\documentclass[11pt,a4paper]{article}
\usepackage[english]{babel}
\usepackage[a4paper,margin=1in]{geometry}
\usepackage[T1]{fontenc}
\usepackage{lmodern}
\usepackage{microtype}
\usepackage{csquotes}
\usepackage{amsmath,amssymb,booktabs,tabularx,array,graphicx}
\usepackage{enumitem}
\usepackage[most]{tcolorbox}
\usepackage{multicol}
\usepackage{wrapfig}
\usepackage[hidelinks]{hyperref}
\usepackage[compact]{titlesec}
\usepackage[font=small,skip=4pt]{caption}
\titlespacing*{\section}{0pt}{8pt}{4pt}
\titlespacing*{\subsection}{0pt}{6pt}{2pt}
\setlength{\textfloatsep}{8pt plus 2pt minus 2pt}
\setlength{\floatsep}{8pt plus 2pt minus 2pt}
\setlength{\intextsep}{8pt plus 2pt minus 2pt}
\setlength{\parskip}{1pt}
\graphicspath{{images/}}

% ---------------------------------------------------------------------------
%  Boxes
% ---------------------------------------------------------------------------
\newtcolorbox{kaobox}[1][]{
  colback=gray!5,
  colframe=gray!55,
  boxrule=0.5pt,
  arc=2pt,
  left=6pt, right=6pt, top=6pt, bottom=6pt,
  #1
}

% ---------------------------------------------------------------------------
%  Figure placeholder: \figph{<width>}{<file name>}
%  Replace by  \includegraphics[width=<width>]{<file name>}  once the image
%  exists in images/.
% ---------------------------------------------------------------------------
\newcommand{\figph}[2]{%
  \fbox{\parbox[c][3.2cm][c]{\dimexpr#1-2\fboxsep-2\fboxrule\relax}{%
    \centering\small\ttfamily [figure: #2]}}}

% ---------------------------------------------------------------------------
%  Shorthand
% ---------------------------------------------------------------------------
\newcommand{\sir}{\mathrm{SIR}}
\newcommand{\gnode}{\mathrm{GN\text{-}ODE}}
\newcommand{\dmp}{\mathrm{DMP}}
\newcommand{\gcn}{\mathrm{GCN}}

\begin{document}
\hypersetup{pageanchor=false}

% ===========================================================================
%  TITLE PAGE
% ===========================================================================
\begin{titlepage}
\centering
\vspace*{0.8cm}

\figph{4cm}{KU\_Logo.png}

\vfill

\rule{\linewidth}{0.5pt}\\[0.7cm]
{\Huge\bfseries Understanding Neural Ordinary Differential Equations\par}
\vspace{0.5cm}
{\LARGE From Simple  Models to Network Epidemics\par}
\vspace{0.7cm}
\rule{\linewidth}{0.5pt}

\vfill

{\large\bfseries Bachelor Seminar}\\[1.6cm]

\begin{tabular}{rl}
\textbf{Presented by} & Veronika Rybak \\[4pt]
\textbf{Supervisor}   & Prof.\ Dr.\ Nadja Ray \\
\end{tabular}

\vfill
{\footnotesize This document accompanies the seminar slides and is intended as a reading and revision handout.}\\[0.3cm]
{\footnotesize\today}
\end{titlepage}

\hypersetup{pageanchor=true}

% ===========================================================================
%  FRONT MATTER
% ===========================================================================
\pagenumbering{roman}
\tableofcontents
\newpage
\pagenumbering{arabic}

\begin{kaobox}[title=Central question of the seminar]
\emph{How can a neural network learn a differential equation?}
A Neural Ordinary Differential Equation does not learn the solution of a system
directly; it learns the rule by which the system changes, and an ODE solver turns
that rule into trajectories. This report builds up this idea on two small
examples that can be checked completely, exponential decay and the SI epidemic
model, and then follows it to a research model for epidemic spreading on contact
networks.
\end{kaobox}

\section*{Abstract}
\addcontentsline{toc}{section}{Abstract}

This report summarizes a bachelor seminar on Neural Ordinary Differential
Equations (Neural ODEs). The first part introduces the necessary theory of
ordinary differential equations, initial value problems, vector fields, and the
Euler method, and explains the Neural ODE as a learned right-hand side of an
ODE. Each concept is illustrated on the exponential decay equation, including a
one-parameter Neural ODE that learns the decay rate.

The second part transfers the same construction to the two-dimensional SI
epidemic model. Own experiments show that a structured Neural ODE recovers the
infection rate from data, that an unstructured network can learn the SI vector
field, and what ``learning the dynamics'' means when the initial condition is
changed. The third part uses this vocabulary to present the Graph Neural ODE
(GN-ODE) model of Kosma et al.\ (2023) for epidemic spreading on contact
networks, where the state consists of one epidemic variable per node and the
interactions are given by the graph. The main conclusion is that Neural ODEs are
most useful when they extend an existing scientific model rather than replace it.


\section{ODE Foundations}

\subsection{Ordinary Differential Equations}

An Ordinary Differential Equation (ODE) describes how a quantity changes over
time instead of directly giving its value. A common first-order ODE has the form
\begin{equation}
  \frac{dy(t)}{dt} = f\bigl(y(t),t\bigr),
  \label{eq:ode}
\end{equation}
where $y(t)$ is the state of the system at time $t$, $dy(t)/dt$ is its rate of
change, and $f$ defines how the system evolves. The unknown in an ODE is a
\emph{function} $y(t)$ rather than a single number, which is the main difference
from an algebraic equation. The word \emph{ordinary} means that the equation
contains derivatives with respect to only one independent variable, here time.

The key idea is that an ODE provides a \emph{local rule of motion}. If we know
the current state and the current time, the equation tells us in which direction
and how fast the state is changing at that moment only; it does not tell us the whole
curve. This idea appears in many applications, including population growth,
radioactive decay, and the motion of a pendulum.

\subsection{A simple scalar ODE: exponential decay}

A well-known example is exponential decay, where a quantity decreases at a rate
proportional to its current value:
\begin{equation}
  \frac{dy}{dt} = -\lambda y, \qquad \lambda > 0 .
  \label{eq:decay}
\end{equation}
Here $f(y,t) = -\lambda y$ does not depend on $t$ explicitly. Its analytical
solution, obtained by separating the variables, is
\begin{equation}
  y(t) = y_0\, e^{-\lambda t}.
  \label{eq:decay-sol}
\end{equation}
Since the rate of change is proportional to the current value, larger values
decrease faster than smaller ones. As a result, the curve becomes flatter as it
approaches zero, as it is demonstrated in Figure 1. The parameter $\lambda$ controls how quickly this happens.

\begin{figure}[h]
  \centering
  \figph{0.6\textwidth}{fig1\_exp\_decay}
  \caption{Exponential decay as a scalar ODE; the decay rate $\lambda$ controls
  how quickly the trajectory falls. Source code is provided in Appendix.}
  \label{fig:decay}
\end{figure}

Exponential decay is used as a running example throughout the first part of the
report: every concept introduced below (initial value problems, vector fields,
numerical solvers, and finally the Neural ODE itself) is first illustrated on
Equation~\eqref{eq:decay}.

\subsection{Initial value problems}

An ODE alone does not determine a single trajectory. For the decay equation, every
function $y(t) = C e^{-\lambda t}$ solves \eqref{eq:decay}, for any constant $C$.
To select one solution, an initial condition is needed. Together,
\begin{equation}
  \frac{dy}{dt} = f(y,t), \qquad y(t_0) = y_0
  \label{eq:ivp}
\end{equation}
form an \emph{initial value problem} (IVP). The ODE defines \emph{how} the system
changes, while the initial condition determines \emph{where} it begins. In the
decay example, $y(0) = y_0$ fixes $C = y_0$ and yields \eqref{eq:decay-sol}.

If $f$ is continuous in $t$ and locally Lipschitz continuous in $y$, then
for every initial condition there exists a unique solution, at least locally
around the initial time (Picard--Lindel\"of theorem). Thus, on an interval
where the solution exists, a given initial state determines one unique
trajectory. This is important for Neural ODEs, where the input is used as the
initial state of an initial value problem.

\begin{figure}[h]
  \centering
  \figph{0.6\textwidth}{fig2\_three\_initial\_values}
  \caption{The same decay rule with three different initial values $y_0$ yields
  three distinct trajectories. Source code is provided in Appendix.}
  \label{fig:three-ivs}
\end{figure}

\subsection{From local derivatives to trajectories}

An ODE tells us how the system changes at its current state. For a scalar
equation such as exponential decay,
\[
    \frac{dy}{dt}=-\lambda y,
\]
the value $-\lambda y$ gives the slope of the trajectory at the current value
of $y$. When $y$ is large, the slope is more negative; as $y$ approaches zero,
the slope becomes smaller. This explains why the decay curve becomes flatter
over time.

For a system with several variables, the state is a vector
$y(t)\in\mathbb{R}^d$, and the function $f$ returns a vector of derivatives.
These derivative vectors describe the direction in which the state moves.
Taken together, they form a \emph{vector field}.

The idea of a vector field is important for Neural ODEs because the neural
network will later play exactly this role: it will learn a function
$f_\theta$ that takes the current state and returns its derivative. A concrete
two-dimensional vector field will be shown later for the SI epidemic model in
Section~\ref{sec:si-field}.

Knowing the derivative at one point is not yet the same as knowing the whole
trajectory. To obtain the trajectory, an ODE solver repeatedly evaluates the
derivative and moves the state forward in time. The simplest example is the
forward Euler method,
\begin{equation}
    y_{n+1}=y_n+h\,f(y_n,t_n),
    \label{eq:euler}
\end{equation}
where $h$ is the step size.

For exponential decay this becomes
\begin{equation}
    y_{n+1}=y_n-h\lambda y_n.
\end{equation}
Starting from $y_0$, Euler's method repeatedly uses the current slope to
approximate the next value. Smaller step sizes usually follow the exact decay
curve more closely, as shown in Figure~\ref{fig:euler}.

\begin{figure}[h]
    \centering
    \figph{0.95\textwidth}{fig4\_euler\_decay}
    \caption{Euler's method approximates the exponential-decay trajectory by
    repeatedly moving in the direction given by the current derivative. A
    smaller step size gives a closer approximation to the exact solution. Source code is provided in Appendix.}
    \label{fig:euler}
\end{figure}

This separation between the two tasks will be important later: the function
$f$ provides the local derivative, while the numerical solver uses these
derivatives to construct the trajectory.

% =====================================================================
\section{From Numerical Solvers to Neural ODEs}

\subsection{Residual networks as motivation}

The historical motivation for Neural ODEs comes from residual networks. A
residual block updates a hidden state according to
\begin{equation}
    h_{k+1}=h_k+F_k(h_k;\theta_k),
\end{equation}
which has the same algebraic form as one Euler step,
\[
    y_{n+1}=y_n+h f(y_n,t_n):
\]
the new state is the current state plus a change computed from that state.

This observation suggests interpreting the layer index $k$ as a discrete time
variable. A residual network can then be viewed as a discrete dynamical system.
Neural ODEs take the corresponding continuous limit: instead of specifying one
update for each discrete layer, a single function $f_\theta$ specifies the
instantaneous rate of change continuously in time. This connection explains the motivation behind Neural ODEs. From this point
on, however, we will focus directly on the ODE viewpoint: the neural network
defines the derivative, and the numerical solver constructs the trajectory.

\subsection{The Neural ODE: core formulation}

A Neural ODE replaces a sequence of discrete updates by a continuous
dynamical system. Instead of moving through states
\[
    h_0 \to h_1 \to \cdots \to h_L,
\]
the state evolves continuously according to
\begin{equation}
    \frac{dh(t)}{dt}
    =
    f_\theta\bigl(h(t),t\bigr),
    \label{eq:node}
\end{equation}
where $f_\theta$ is a neural network with trainable parameters $\theta$.

The role of $f_\theta$ is the same as the function $f$ in an ordinary
differential equation: given the current state, it returns the instantaneous
rate of change. Starting from an initial state $h(0)$, an ODE solver repeatedly
evaluates $f_\theta$ and uses these derivatives to construct the trajectory
$h(t)$.

Formally, the state at a later time $T$ satisfies
\begin{equation}
    h(T)
    =
    h(0)
    +
    \int_0^T
    f_\theta\bigl(h(t),t\bigr)\,dt .
    \label{eq:odesolve}
\end{equation}

The important difference from an ordinary neural network is therefore that
$f_\theta$ does not directly compute the final output. It defines how the
state changes locally, while the ODE solver combines these local changes to
produce the trajectory.

\paragraph{A simple trainable example: exponential decay.}

The exponential-decay example can be used to show the learning idea in the
simplest possible setting. Suppose the true system is
\[
    \frac{dy}{dt}=-\lambda y,
\]
but the decay rate $\lambda$ is unknown. We replace it by a trainable
parameter $\theta$ and use
\begin{equation}
    f_\theta(y)=-\theta y.
    \label{eq:toy-field}
\end{equation}

The resulting model is
\[
    \frac{dy}{dt}=-\theta y,
    \qquad
    y(0)=y_0.
\]
For a given value of $\theta$, the ODE solver produces a trajectory. This
trajectory can be compared with observed data, and $\theta$ can then be
adjusted to reduce the difference. If the data were generated using the true
decay rate $\lambda$, training should recover a value of $\theta$ close to
$\lambda$.

This example does not yet learn an unknown functional form: the structure
$-\theta y$ is already built into the model, and only one parameter is
learned. Its purpose is to make the Neural ODE training idea transparent
before moving to a neural network that learns the derivative function itself.

\subsection{What the network learns}

A Neural ODE learns the function $f_\theta$ that describes how the current
state changes. The ODE solver then uses this function to construct the
trajectory.

In the exponential-decay example,
\[
    f_\theta(y)=-\theta y,
\]
so training changes the parameter $\theta$. Once $\theta$ has been learned,
the same differential equation can be solved from different initial values
$y_0$. The initial value determines where the trajectory starts, while
$f_\theta$ determines how it evolves.

For this simple example, the functional form $-\theta y$ is already known and
only the decay rate is learned. Later, in the SI example, $f_\theta$ will also
be represented by a neural network, so that more of the derivative function
can be learned from data.

\begin{figure}[h]
    \centering
    \figph{0.6\textwidth}{fig7\_chen\_resnet\_vs\_ode}
    \caption{Residual network (left) and Neural ODE (right). A residual network
    performs a sequence of discrete updates, whereas a Neural ODE describes the
    state continuously through a differential equation. Original figure:
    Chen et al.\ (2018), Fig.~1.}
    \label{fig:chen}
\end{figure}

\subsection{Training a Neural ODE}

Training changes the parameters $\theta$ of the derivative function
$f_\theta$. The basic procedure is:

\begin{enumerate}
    \item[(1)] \textbf{Forward solve:}
    starting from the initial state, the ODE solver uses the current
    $f_\theta$ to produce a predicted trajectory.

    \item[(2)] \textbf{Loss:}
    the predicted states are compared with the observed data. For example,
    for observations $y^*(t_k)$ we can use
    \[
        L(\theta)
        =
        \sum_k
        \left(
            y_\theta(t_k)-y^*(t_k)
        \right)^2 .
    \]

    \item[(3)] \textbf{Gradient computation:}
    the loss is differentiated with respect to the parameters $\theta$.

    \item[(4)] \textbf{Parameter update:}
    the parameters are changed in a direction that reduces the loss,
    \[
        \theta
        \leftarrow
        \theta-\eta\frac{\partial L}{\partial\theta},
    \]
    where $\eta$ is the learning rate.
\end{enumerate}

The important point is that training does not directly modify the trajectory.
It modifies $f_\theta$. When the ODE is solved again with the updated
$f_\theta$, a different trajectory is produced.

\begin{figure}[h]
    \centering
    \figph{0.95\textwidth}{fig8\_training\_loop}
    \caption{Training a Neural ODE: solve the ODE, compare the predicted
    trajectory with the data, compute the gradient, and update the parameters
    of $f_\theta$.}
    \label{fig:training}
\end{figure}

\paragraph{Differentiating through the solver.}

To update the parameters $\theta$, the loss must be differentiated with
respect to them. Since the predicted trajectory is produced by the ODE solver,
this means differentiating through the solver steps. For the simple Euler-based
experiments in this report, this can be done directly using the chain rule.

\begin{figure}[h]
    \centering
    \figph{0.95\textwidth}{fig\_toy\_neural\_ode}
    \caption{Training the model $dy/dt=-\theta y$. Left: the trajectory before
    and after training compared with the reference trajectory. Right: the
    learned parameter $\theta$ approaches the true decay rate $\lambda$.}
    \label{fig:toy}
\end{figure}


\section{From Exponential Decay to a Two-Dimensional System: the SI Model}
\label{sec:si}
\subsection{The SI equations}

So far the state of the system was a single number $y(t)$. Epidemic models
describe how a population moves between health states, so several quantities
change at the same time. The simplest such model is the \emph{SI model}, in which
individuals are either susceptible ($S$) or infected ($I$) and, once infected,
stay infected. We use population \emph{fractions}, so that
\begin{equation}
  S(t) + I(t) = 1 .
  \label{eq:si-conservation}
\end{equation}
The dynamics are
\begin{equation}
  \frac{dS}{dt} = -\beta\, S I, \qquad \frac{dI}{dt} = \beta\, S I,
  \label{eq:si}
\end{equation}
with a single parameter $\beta>0$, the infection (or contact) rate.

\paragraph{Meaning of the terms.}
The product $SI$ measures how often susceptible and infected individuals meet:
if either group is small, meetings are rare. Each meeting transmits the disease
at rate $\beta$. The term $-\beta SI$ removes these individuals from the
susceptible group, and the same amount $+\beta SI$ enters the infected group.
Adding the two equations gives $d(S+I)/dt = 0$, which is why
\eqref{eq:si-conservation} holds for all times, not only at $t=0$.

\paragraph{Connection to the scalar example.}
In Section~1 the ODE was $dy/dt = -\lambda y$ with one state variable. The SI
model has the same structure with two state variables. Writing the state as a
vector $\mathbf{x} = (S,I)^\top$, the two equations \eqref{eq:si} become one
vector-valued ODE,
\begin{equation}
  \frac{d}{dt}\begin{pmatrix} S \\ I \end{pmatrix}
  = \begin{pmatrix} -\beta SI \\ \beta SI \end{pmatrix}
  = f(\mathbf{x}),
  \qquad \text{i.e.}\qquad \dot{\mathbf{x}} = f(\mathbf{x}),
  \label{eq:si-vector}
\end{equation}
which is exactly the form \eqref{eq:ode} with $d=2$. Everything from
Section~1 carries over: an initial condition $(S(0),I(0))$ selects one
trajectory, $f$ is a vector field on the $(S,I)$-plane, and the Euler method
\eqref{eq:euler} applies componentwise. One difference is that
$f(\mathbf{x})$ is now nonlinear in the state (because of the product $SI$), so
there is no simple closed formula to compare with; we rely on the numerical
solver.

\paragraph{A first simulation.}
Figure~\ref{fig:si-sim} shows the Euler solution of \eqref{eq:si} with
$\beta=1$, $S(0)=0.99$, $I(0)=0.01$, step size $h=0.01$. Three phases are
visible. At the beginning $I$ is small, so the product $SI$ is small and the
infected fraction grows slowly. As $I$ increases while $S$ is still large, $SI$
becomes large and the growth is fastest (around $t\approx 4$--$5$, where
$S\approx I\approx 0.5$ and $SI$ takes its maximum value $0.25$). Finally $S$
becomes small, $SI$ shrinks again, and $I$ approaches $1$: without recovery,
eventually everyone is infected. The dotted line confirms that $S+I$ stays equal
to one throughout.

\begin{figure}[h]
  \centering
  \figph{0.65\textwidth}{fig\_si\_simulation}
  \caption{Numerical solution of the SI model \eqref{eq:si} with $\beta=1$,
  $S(0)=0.99$, $I(0)=0.01$. The infected fraction follows an S-shaped curve; the
  sum $S+I$ is conserved.}
  \label{fig:si-sim}
\end{figure}

\subsection{The SI model as a vector field}
\label{sec:si-field}

Section~1.4 called $f$ a vector field. For the SI model this can be drawn:
in the plane with coordinates $S$ and $I$, every point $(S,I)$ carries the
vector $f(S,I)$ of \eqref{eq:si-vector}. Figure~\ref{fig:si-field} shows
these vectors on a grid together with the trajectory of Figure~\ref{fig:si-sim}.
This picture is needed once, because it is exactly what the neural network in
Section~\ref{sec:si-node} will have to learn.

A single point makes the picture concrete. At $(S,I) = (0.8,\,0.2)$ with
$\beta=1$,
\[
  f(0.8,0.2) = (-0.8\cdot 0.2,\; 0.8\cdot 0.2) = (-0.16,\; 0.16),
\]
so at that moment $S$ decreases at rate $0.16$ and $I$ increases at exactly the
same rate: the arrow points up and to the left at $45^\circ$. Because the two
components are always equal and opposite, \emph{every} arrow in the plane has
this direction; only the length $\sqrt{2}\,\beta SI$ varies. Arrows are longest
near $S=I=0.5$ and vanish on the axes $S=0$ and $I=0$, where no meetings occur.

Since the arrows are parallel to the line $S+I=1$, a trajectory starting on
this line stays on it, which is the conservation law \eqref{eq:si-conservation}
seen geometrically; and since the arrows are longest in the middle, the
trajectory moves fastest there, which is the S-shape of Figure~\ref{fig:si-sim}
seen from a different angle.

\begin{figure}[h]
  \centering
  \figph{0.55\textwidth}{fig\_si\_phase\_plane}
  \caption{The SI vector field for $\beta=1$ and the trajectory starting at
  $(0.99,0.01)$. All arrows point in the direction $(-1,1)$; their length is
  $\sqrt{2}\,SI$. The trajectory moves along the line $S+I=1$ from the lower right
  to the upper left, fastest in the middle.}
  \label{fig:si-field}
\end{figure}

\section{Turning the SI Model into a Neural ODE}
\label{sec:si-node}

For exponential decay we replaced the known right-hand side $-\lambda y$ by a
learned function $f_\theta(y)$ (Section~2.2). We now do the same for the SI
model. Starting from
\[
  \dot{\mathbf{x}} = \begin{pmatrix} -\beta SI \\ \beta SI \end{pmatrix},
\]
a Neural ODE replaces some or all of the right-hand side by a learned function
of the state,
\begin{equation}
  \dot{\mathbf{x}} = f_\theta(\mathbf{x}), \qquad \mathbf{x} = (S,I)^\top,
  \qquad f_\theta : \mathbb{R}^2 \to \mathbb{R}^2 .
  \label{eq:si-node}
\end{equation}
Nothing conceptually new has happened compared with exponential decay. The state
has changed from one number $y$ to two numbers $S$ and $I$, and the learned
derivative therefore has two outputs instead of one.

\paragraph{Architecture.}
The complete model is the following pipeline:
\[
  (S,I)\;\longrightarrow\;
  \boxed{\text{small neural network } f_\theta}\;\longrightarrow\;
  \Bigl(\tfrac{dS}{dt},\tfrac{dI}{dt}\Bigr)\;\longrightarrow\;
  \boxed{\text{ODE solver}}\;\longrightarrow\;
  \bigl(S(t+\Delta t),\, I(t+\Delta t)\bigr).
\]
The network never sees a time series and never outputs $S(t)$ or $I(t)$
directly. It receives the \emph{current state} and returns the
\emph{instantaneous rate of change} of that state. The solver then advances the
state by repeatedly asking the network for the derivative, exactly as the Euler
method repeatedly evaluated $-\lambda y$ in Section~1.4. This division of labour
is the central point of the report: the network learns the vector field, the
solver turns it into trajectories.

\paragraph{Two ways to choose $f_\theta$.}
How much of the right-hand side is learned is a modelling choice.
\begin{itemize}
  \item \emph{Structured (grey-box):} keep the SI form and learn only the
  parameter, $f_\theta(S,I) = (-\theta SI,\ \theta SI)^\top$. This is the exact
  analogue of $f_\theta(y) = -\theta y$ in Section~2.2.
  \item \emph{Unstructured (black-box):} let a small multilayer perceptron with
  input $(S,I)$ and output $(\dot S,\dot I)$ learn the whole vector field, without
  being told the SI form.
\end{itemize}
Both are trained with the four steps of Section~2.4: solve forward from the
initial condition, compare the resulting $S(t_k), I(t_k)$ with data at the
observation times $t_k$, compute the gradient, update $\theta$.

\section{Neural ODE Experiments on the SI Model}
\label{sec:si-experiments}

All data in this section are generated from the SI model \eqref{eq:si} with
$\beta_{\text{true}} = 1$ and solved with the Euler method, $h = 0.01$, on
$t\in[0,15]$. Because the data-generating equation is known, every result can be
checked against the truth.

\subsection{Experiment 1: recovering a known parameter}
\label{sec:exp1}

We pretend $\beta$ is unknown and use the structured model
$f_\theta(S,I) = (-\theta SI,\ \theta SI)^\top$ with a single learnable parameter
$\theta$. The training data are the values $I^*(t_k)$ of the true solution at
$t_k = 1,2,\dots,15$, starting from $(S_0,I_0) = (0.99,0.01)$. The loss is the
squared error summed over the observation times,
\[
  L(\theta) = \sum_{k} \bigl(I_\theta(t_k) - I^*(t_k)\bigr)^2,
\]
where $I_\theta(t_k)$ is the second component of the solver output for the
current $\theta$. The gradient $\partial L/\partial\theta$ is computed by
differentiating through the Euler steps, and $\theta$ is updated by gradient
descent with learning rate $\eta = 0.02$, starting from $\theta_0 = 0.3$.

Figure~\ref{fig:si-beta} shows the result. Before training the model's epidemic
is far too slow (dashed red). After 60 iterations $\theta = 1.000$, and the
predicted curve lies on top of the data. The right panel shows that $\theta$
moves quickly in the first few iterations, when the mismatch is large, and then
settles; the same qualitative behaviour as for the decay rate in
Figure~\ref{fig:toy}.

It is important to state precisely what has been learned. In this experiment the
Neural ODE does not discover an unknown functional form: the SI structure is
supplied by the model, and training identifies the unknown parameter $\beta$.
The value of the experiment is that it shows the training loop of Section~2.4
working on a two-dimensional, nonlinear system, with an answer we can verify.

\begin{figure}[h]
  \centering
  \figph{0.95\textwidth}{fig\_si\_beta\_recovery}
  \caption{Experiment 1. Left: infected fraction from the true model (solid),
  the observations used for training (dots), and the structured Neural ODE
  before and after training (dashed). Right: the learned parameter $\theta$
  converges to $\beta_{\text{true}}=1$.}
  \label{fig:si-beta}
\end{figure}

\subsection{Experiment 2: learning the whole SI vector field}
\label{sec:exp2}

% --- TO RUN. Specification: ---
% f_theta = MLP 2 -> 16 -> 16 -> 2 (tanh), Euler solver h=0.01 inside the loss,
% training data: full trajectories S(t),I(t) from (0.99,0.01), (0.95,0.05),
% (0.90,0.10); loss = MSE over all time points; Adam, ~2000 steps.
% Figure fig_si_learned_field.png: left true field f(S,I), right learned f_theta(S,I),
% same grid and arrow scale; overlay the training trajectory on both.

The structured model could only ever produce SI-shaped curves. We now remove
this prior knowledge and let a small multilayer perceptron
$f_\theta:\mathbb{R}^2\to\mathbb{R}^2$ (input $(S,I)$, output $(\dot S,\dot I)$,
two hidden layers) learn the vector field from trajectories generated by
\eqref{eq:si}. The training loop is unchanged; only the definition of $f_\theta$
differs.

Figure~\ref{fig:si-learned-field} compares the true vector field $f(S,I)$ with
the learned $f_\theta(S,I)$ on the same grid as Figure~\ref{fig:si-field}.
[Describe: where the arrows agree (along and near the training trajectories),
where they do not (regions far from $S+I=1$, which the training data never
visited), and that the learned field need not conserve $S+I$ exactly.] With this
figure the statement ``the neural network learns the vector field'' is no longer
a general claim; it is something that was trained, plotted, and compared with the
truth.

\begin{figure}[h]
  \centering
  \figph{0.9\textwidth}{fig\_si\_learned\_field}
  \caption{Experiment 2. Left: true SI vector field. Right: vector field learned
  by the unstructured Neural ODE. Training trajectories overlaid.}
  \label{fig:si-learned-field}
\end{figure}

\subsection{Experiment 3: an unseen initial condition}
\label{sec:exp3}

% --- TO RUN. Use the trained MLP from Experiment 2. Test initial condition
% (0.93,0.07), not in the training set. Figure fig_si_unseen_ic.png: true I(t)
% vs predicted I(t) for the test IC; optionally also one training IC for reference.

The model of Experiment~2 was trained on trajectories starting from
$I_0 \in \{0.01, 0.05, 0.10\}$. We now start the solver with the learned
$f_\theta$ from $I_0 = 0.07$, an initial condition that was never used during
training, and compare with the true solution (Figure~\ref{fig:si-unseen}).

[Describe the agreement.] This is the point where the IVP discussion of
Section~1.3 pays off. Changing the initial condition does not require learning a
new trajectory. The solver integrates the same learned differential equation
from a different starting point, just as Figure~\ref{fig:three-ivs} showed three
decay curves produced by one rule. The experiment also shows the limit of this
statement: the new initial condition lies inside the region of the $(S,I)$-plane
covered by the training data. Starting from a point far outside that region
would query $f_\theta$ where it was never trained, and the trajectory would not
be reliable.

\begin{figure}[h]
  \centering
  \figph{0.9\textwidth}{fig\_si\_unseen\_ic}
  \caption{Experiment 3. True and predicted $I(t)$ for the unseen initial
  condition $I_0=0.07$.}
  \label{fig:si-unseen}
\end{figure}

\section{What the Simple Examples Teach}
\label{sec:lessons}

Before turning to the network model of Kosma et al., we collect what the
experiments have shown. During prediction, a Neural ODE is the chain
\[
  \boxed{\text{state}} \rightarrow \boxed{f_\theta} \rightarrow
  \boxed{\text{derivative}} \rightarrow \boxed{\text{ODE solver}} \rightarrow
  \boxed{\text{trajectory}},
\]
and during training the chain is closed by
\[
  \boxed{\text{trajectory}} \rightarrow \boxed{\text{loss}} \rightarrow
  \boxed{\text{gradient}} \rightarrow \boxed{\theta\ \text{update}}.
\]
Three points carry over to any Neural ODE, however large:
\begin{enumerate}
  \item The network models the \emph{derivative}, not the solution. The solution
  is produced by the solver (Sections~2.3 and~\ref{sec:si-node}).
  \item Moving from one state variable to two changed nothing in the method; only
  the dimension of the input and output of $f_\theta$ changed
  (Section~\ref{sec:si-node}). Moving to many variables, as in the next part,
  is the same step repeated.
  \item How much structure is built into $f_\theta$ is a choice. A structured
  $f_\theta$ (Experiment~1) learns little but is interpretable and needs almost
  no data; an unstructured $f_\theta$ (Experiment~2) is flexible but reliable
  only where data exist (Experiment~3). The model of the next part sits in
  between: it keeps the epidemic structure and learns the parts that are unknown.
\end{enumerate}

% ===========================================================================
%  % ===========================================================================
\section{From the SI Model to Epidemic Spreading on Networks}
\label{sec:networks}
% ===========================================================================

The SI example was deliberately simple: the complete state consisted of only
two numbers, $S(t)$ and $I(t)$, and every individual was treated as part of one
well-mixed population. The application studied by Kosma et al.\ (2023) retains
the same basic idea---describe local rates of change by an ODE---but enlarges
the state in two steps. First, recovery is added. Second, instead of describing
only population fractions, the epidemic state is represented separately for
the nodes of a contact network.


\subsection{Adding recovery: from SI to SIR}

The SI model assumes that an infected individual remains infected forever. The
SIR model adds a recovered state $R$. For population fractions,
\[
    S(t)+I(t)+R(t)=1,
\]
and the dynamics are
\begin{align}
    \frac{dS}{dt} &= -\beta SI, \\
    \frac{dI}{dt} &=  \beta SI-\gamma I, \\
    \frac{dR}{dt} &=  \gamma I .
\end{align}
The new parameter $\gamma>0$ is the recovery rate. The infection term is the
same one already studied in the SI model. The only new process is
\[
    I \longrightarrow R,
\]
which removes infected individuals at rate $\gamma I$.

Thus the move from SI to SIR does not introduce a new modelling principle. It
only enlarges the state from
\[
    \mathbf{x}=(S,I)^\top
\]
to
\[
    \mathbf{x}=(S,I,R)^\top
\]
and changes the vector field accordingly.


\subsection{From one population to one state per node}

The well-mixed equations still assume that every susceptible individual can
interact equally with every infected individual. To model a contact structure,
let
\[
    G=(V,E)
\]
be a graph with $n=|V|$ nodes and adjacency matrix
$A\in\mathbb{R}^{n\times n}$. The entry $A_{ij}$ describes whether node $i$
can interact with node $j$.

For every node $i$, define
\[
    s_i(t),\qquad i_i(t),\qquad r_i(t)
\]
as the probabilities of being susceptible, infectious, and recovered. They
satisfy
\[
    s_i+i_i+r_i=1.
\]
Collecting all nodes gives the vectors
\[
    \mathbf{s},\mathbf{i},\mathbf{r}\in\mathbb{R}^n.
\]

The complete epidemic state can therefore be written as the matrix
\begin{equation}
    X(t)=
    \begin{pmatrix}
        s_1(t) & i_1(t) & r_1(t)\\
        s_2(t) & i_2(t) & r_2(t)\\
        \vdots & \vdots & \vdots\\
        s_n(t) & i_n(t) & r_n(t)
    \end{pmatrix}
    \in\mathbb{R}^{n\times 3}.
    \label{eq:network-state-matrix}
\end{equation}
Each row describes one node and each column one epidemic state.

This is the network analogue of the vector
$(S,I)^\top$ used in Section~\ref{sec:si}. The main difference is dimension:
instead of evolving two population-level variables, the model now evolves
state information for every node.


\subsection{A vector form of the approximate network dynamics}

Under a simple independence approximation, the node-level equations can be
written compactly as
\begin{equation}
    \frac{d}{dt}
    \begin{pmatrix}
        \mathbf{s}\\
        \mathbf{i}\\
        \mathbf{r}
    \end{pmatrix}
    =
    \begin{pmatrix}
        -\beta\,\mathbf{s}\odot(A\mathbf{i})\\[3pt]
         \beta\,\mathbf{s}\odot(A\mathbf{i})-\gamma\mathbf{i}\\[3pt]
         \gamma\mathbf{i}
    \end{pmatrix},
    \label{eq:network-sir-vector}
\end{equation}
where $\odot$ denotes elementwise multiplication.

Equation~\eqref{eq:network-sir-vector} is the direct network counterpart of
Equation~\eqref{eq:si-vector}. The vector $A\mathbf{i}$ has one entry per node.
Its $i$th entry,
\[
    (A\mathbf{i})_i=\sum_j A_{ij}i_j,
\]
summarizes the infectious pressure coming from the neighbours of node $i$.
Multiplying this quantity by $s_i$ gives the rate at which node $i$ loses
susceptibility.

This representation is useful because it shows that moving from the SI model
to a network has not changed the Neural ODE principle. We still have
\[
    \dot{\mathbf{x}}=f(\mathbf{x}),
\]
but $\mathbf{x}$ is now a much larger state and the vector field also depends
on the graph structure $A$.


\subsection{Why the simple network equations are only an approximation}

There is an important difficulty hidden in
Equation~\eqref{eq:network-sir-vector}. In the exact stochastic network model,
the infection probability of node $i$ depends on a joint event: node $i$ must
be susceptible while one of its neighbours $j$ is infected. For example,
\begin{equation}
    \frac{dP(X_i=S)}{dt}
    =
    -\beta\sum_j
    A_{ij}\,
    P(X_i=S,\;X_j=I).
    \label{eq:exact-node}
\end{equation}
The individual probabilities $P(X_i=S)$ and $P(X_j=I)$ are therefore not
enough to evaluate the derivative.

A simple closure replaces the joint probability by
\[
    P(X_i=S,\;X_j=I)
    \approx
    P(X_i=S)P(X_j=I).
\]
This produces the compact equations above, but it ignores correlations between
neighbouring nodes.

The exact problem also becomes large very quickly. Since each of $n$ nodes can
be in one of the three states $S$, $I$, or $R$, the complete stochastic
network has $3^n$ possible global configurations. Rather than explicitly
tracking this exponentially large state space, practical models introduce
approximations. GN-ODE can be understood as one attempt to retain the useful
SIR structure while adding learnable components that compensate for some of
the information lost by the simplified representation.


% ===========================================================================
\section{GN-ODE as an Extension of the Simple Neural ODE}
\label{sec:gnode}
% ===========================================================================

The examples in Sections~\ref{sec:si-node}--\ref{sec:si-experiments} make it
possible to describe the Graph Neural ODE without introducing a completely new
set of concepts.

For the SI experiment, the model had the form
\[
    \dot{\mathbf{x}}=f_\theta(\mathbf{x}),
    \qquad
    \mathbf{x}=(S,I)^\top.
\]
For a network, the same idea becomes schematically
\[
    \dot{X}=f_\theta(X,A),
\]
where $X$ contains the node states and $A$ supplies the information about which
nodes interact.


\subsection{Encode, evolve, decode}

The GN-ODE model of Kosma et al.\ (2023) follows an
\emph{encode--evolve--decode} structure.

First, the observed node states are transformed into hidden representations by
a trainable encoder. Instead of evolving the binary input states directly, the
model works with hidden susceptible, infectious, and recovered
representations.

Second, these hidden variables evolve inside an ODE solver. Their dynamics
retain the structure of the SIR equations. In simplified notation,
\begin{align}
    \frac{dS}{dt}
        &= -\beta\,(A I_h)\odot S_h,\\
    \frac{dI}{dt}
        &=  \beta\,(A I_h)\odot S_h-\gamma I_h,\\
    \frac{dR}{dt}
        &=  \gamma I_h.
    \label{eq:gnode-core}
\end{align}
Here $S_h$ and $I_h$ are learned hidden representations, while multiplication
by $A$ aggregates information from neighbouring nodes.

Finally, a trainable decoder maps the evolved hidden states back to
probabilities of the observable states $S$, $I$, and $R$.

The important point is that the ODE solver has exactly the same role as in the
simple examples. It does not learn anything. It repeatedly evaluates the
right-hand side and integrates the state forward in time. The learnable
components determine the representation and therefore the vector field that
the solver follows.


\subsection{Relation to the previous examples}

Table~\ref{tab:progression} summarizes the progression of the report.

\begin{table}[h]
\centering
\small
\begin{tabularx}{\textwidth}{lXXX}
\toprule
Model & State & Right-hand side & What is learned\\
\midrule
Exponential decay
&
$y$
&
$-\theta y$
&
decay rate $\theta$
\\

Structured SI Neural ODE
&
$(S,I)$
&
$(-\theta SI,\theta SI)$
&
infection rate $\theta$
\\

Black-box SI Neural ODE
&
$(S,I)$
&
MLP $f_\theta(S,I)$
&
complete local vector field
\\

GN-ODE
&
node-level hidden states
&
SIR-shaped graph dynamics
&
hidden representations and mappings
\\
\bottomrule
\end{tabularx}
\caption{The same Neural ODE idea at increasing levels of complexity.}
\label{tab:progression}
\end{table}

The table also clarifies why the simple experiments are useful. GN-ODE is not
a fundamentally different mechanism. It is a larger, graph-structured version
of the same construction:
\[
    \text{state}
    \rightarrow
    \text{learned local dynamics}
    \rightarrow
    \text{ODE solver}
    \rightarrow
    \text{trajectory}.
\]


% ===========================================================================
\section{Published Evaluation of GN-ODE}
\label{sec:published-results}
% ===========================================================================

The experiments in Section~\ref{sec:si-experiments} are small simulations
performed specifically for this report. This section is different: it
summarizes the experiments reported by Kosma et al.\ (2023) and should
therefore be interpreted as results from the literature rather than as my own
experiments.

The authors evaluate GN-ODE on several social and communication networks of
different sizes. Epidemic trajectories are generated for different infection
and recovery parameters, and Monte Carlo simulations are used as the reference
solution. GN-ODE is compared with both learning-based and classical epidemic
prediction methods.


\subsection{Prediction on the same network}

In the first setting, training and test trajectories come from the same graph,
but the epidemic parameters vary between simulations. This asks whether the
model has learned dynamics that remain useful when the infection and recovery
rates change.

Kosma et al.\ report that GN-ODE achieves prediction errors comparable to or
lower than the learning-based baselines on the networks considered. Classical
Dynamic Message Passing remains highly accurate in several settings, but its
computational cost becomes increasingly important on large graphs.

The relevant interpretation for this report is not that one method always
dominates the others. Rather, GN-ODE attempts to occupy an intermediate
position: it keeps explicit epidemic structure, gains flexibility from learned
representations, and can be evaluated without performing a large number of
Monte Carlo simulations for every new prediction.


\subsection{Transfer to larger unseen networks}

A second experiment is particularly relevant to the idea of learning a
dynamical rule rather than memorizing one trajectory. The model is trained on
smaller graphs and then evaluated on larger networks that were not present
during training.

Kosma et al.\ report that GN-ODE transfers more successfully than the standard
graph-neural baseline in this setting. Dynamic Message Passing remains a strong
accuracy reference, but it is an algorithm that must be executed again on each
new graph rather than a trained model whose parameters can be reused.

This experiment is the network-scale analogue of changing the initial
condition in Section~\ref{sec:exp3}. In both cases, the objective is to reuse
the same learned local rule in a situation that was not presented exactly
during training. The network experiment is substantially harder, however,
because not only the initial state but also the number of nodes and the graph
structure can change.


% ===========================================================================
\section{Critical Analysis}
\label{sec:critical}
% ===========================================================================

The simple experiments and the GN-ODE application illustrate both the appeal
and the limitations of Neural ODEs.

\subsection{What the simple experiments establish}

The decay and SI experiments make three aspects of the method directly
observable. First, the object being learned is a derivative function rather
than a complete solution curve. Second, the ODE solver and the neural network
have separate roles: the network supplies local rates of change and the solver
integrates them. Third, prior scientific knowledge can be incorporated at
different levels. Learning only $\beta$ in the structured SI model assumes
almost the complete equation, whereas the black-box SI experiment places much
less structure in the model.

These examples are intentionally controlled. The training data are generated
from equations that are already known and contain no observational noise.
Their purpose is therefore not to demonstrate an advantage over classical
epidemic modelling. Their purpose is to make the Neural ODE mechanism
verifiable.


\subsection{A limitation visible already in the SI experiment}

The unstructured SI experiment also illustrates an important limitation.
Because normalized SI trajectories satisfy
\[
    S+I=1,
\]
all physically valid training data lie on a one-dimensional subset of the
two-dimensional $(S,I)$ plane. The loss therefore constrains
$f_\theta(S,I)$ mainly on or near this line. Its behaviour far away from the
observed trajectories is not identified by the data.

For the same reason, testing the model at $I_0=0.07$ is useful as an
interpolation test, but it should not be interpreted as strong evidence of
general out-of-distribution generalization. The new state still lies on the
same invariant manifold and between initial conditions seen during training.

This observation is important for larger Neural ODEs as well: a learned vector
field is only supported empirically in regions of state space that are
sufficiently constrained by the training trajectories.


\subsection{Strengths and limitations of GN-ODE}

The main strength of GN-ODE is the combination of explicit model structure and
learning. The SIR mechanism and graph interactions are not discarded, so the
model begins with useful scientific information instead of learning epidemic
dynamics entirely from scratch. At the same time, trainable hidden
representations give the model more flexibility than a fixed mean-field
approximation.

There are also costs. The model depends on repeated numerical solution of an
ODE, so computational cost and numerical error depend on the chosen solver and
step size or tolerance. In addition, although the SIR structure remains
visible, the hidden representations produced by the learned transformations
are less directly interpretable than the original probabilities
$s_i$, $i_i$, and $r_i$.

Most importantly, the presence of an ODE does not by itself guarantee correct
extrapolation. The vector field is still learned from finite training data.
The same warning seen in the SI experiment therefore remains valid on a much
larger scale: predictions are most credible when the training data constrain
the regions of state space and graph structures that the model will encounter.


% ===========================================================================
\section{Summary and Conclusion}
% ===========================================================================

This report approached Neural Ordinary Differential Equations from a sequence
of models whose behaviour could first be understood directly.

The starting point was exponential decay,
\[
    \dot y=-\lambda y.
\]
This example separated three objects that remain distinct throughout the
report: the current state, the derivative rule, and the trajectory produced by
solving the ODE. Replacing the known decay rate by a trainable parameter then
showed, in the simplest possible setting, what learning an ODE means.

The same construction was next applied to the SI epidemic model. The state
became two-dimensional,
\[
    \mathbf{x}=(S,I)^\top,
\]
but the Neural ODE mechanism did not change. A structured model learned an
unknown infection rate, while an unstructured model treated the complete
vector field as the learnable object. The SI example also made the geometric
meaning of a vector field visible and showed why predictions from unseen
initial conditions are obtained by integrating the same learned dynamics from
a new starting state.

Finally, the same viewpoint was used to interpret GN-ODE. The state is enlarged
from a few population variables to representations associated with all nodes
of a contact graph, and the adjacency matrix determines which nodes can
interact. Nevertheless, the central computation remains
\[
    \text{state}
    \rightarrow
    f_\theta
    \rightarrow
    \text{derivative}
    \rightarrow
    \text{ODE solver}
    \rightarrow
    \text{trajectory}.
\]

The main lesson is therefore not that Neural ODEs replace differential-equation
models. Their usefulness is clearest when the two approaches are combined:
known mechanisms can determine part of the vector field, while learnable
components represent effects that are difficult to specify analytically. The
simple decay and SI experiments make this principle explicit; GN-ODE shows how
the same principle can be extended to a substantially larger scientific
problem.


% ===========================================================================
\section{References}
% ===========================================================================

{\small
\begin{description}[leftmargin=*,itemsep=2pt]

\item[Chen et al.\ (2018)]
Chen, R.\ T.\ Q., Rubanova, Y., Bettencourt, J., \& Duvenaud, D.\ K.
\emph{Neural Ordinary Differential Equations}.
Advances in Neural Information Processing Systems 31.

\item[Kosma et al.\ (2023)]
Kosma, C., Nikolentzos, G., Panagopoulos, G., Steyaert, J.-M., \&
Vazirgiannis, M.
\emph{Neural Ordinary Differential Equations for Modeling Epidemic Spreading}.
Transactions on Machine Learning Research.

\end{description}
}

\end{document}